\chapter{2005年北京卷（秋文）}
\section{单选题}
\begin{problem}
设集合 \(U=\R\)，集合 \(M=\{x\mid x>1\}\)，\(P=\{x\mid x^2>1\}\)，则下列关系中正确的是（\quad）

\choices
    {\(M=P\)}
    {\(P\not\subseteq M\)}
    {\(M\not\subseteq P\)}
    {\(M\cap P=\R\)}
\begin{answer}
B
\end{answer}
\end{problem}

\begin{solution}
由 \(x^2>1\) 得 \(x<-1\) 或 \(x>1\)，所以 \(P=(-\infty,-1)\cup(1,+\infty)\)，而 \(M=(1,+\infty)\). 因此 \(M\subseteq P\)，但 \(P\not\subseteq M\)，故选 B.
\end{solution}

\begin{problem}
为了得到函数 \(y=2^{x-3}-1\) 的图象，只需把函数 \(y=2^x\) 的图象上所有的点（\quad）

\choices
    {向右平移 \(3\) 个单位长度，再向下平移 \(1\) 个单位长度}
    {向左平移 \(3\) 个单位长度，再向下平移 \(1\) 个单位长度}
    {向右平移 \(3\) 个单位长度，再向上平移 \(1\) 个单位长度}
    {向左平移 \(3\) 个单位长度，再向上平移 \(1\) 个单位长度}
\begin{answer}
A
\end{answer}
\end{problem}

\begin{solution}
将 \(y=2^x\) 中的 \(x\) 换成 \(x-3\)，所得图象向右平移 \(3\) 个单位长度；再减去 \(1\)，所得图象向下平移 \(1\) 个单位长度，故选 A.
\end{solution}

\begin{problem}
“\(m=\frac{1}{2}\)”是“直线 \((m+2)x+3my+1=0\) 与直线 \((m-2)x+(m+2)y-3=0\) 相互垂直”的（\quad）

\choices
    {充分必要条件}
    {充分而不必要条件}
    {必要而不充分条件}
    {既不充分也不必要条件}
\begin{answer}
B
\end{answer}
\end{problem}

\begin{solution}
两直线相互垂直的条件为 \((m+2)(m-2)+3m(m+2)=0\)，即 \(2(2m-1)(m+2)=0\)，所以 \(m=\frac12\) 或 \(m=-2\). 当 \(m=\frac12\) 时两直线确实垂直；当 \(m=-2\) 时两直线分别为 \(y=\frac16\) 和 \(x=-\frac34\)，也相互垂直. 因此“\(m=\frac12\)”是充分而不必要条件，故选 B.
\end{solution}

\begin{problem}
若 \(|\bs{a}|=1\)，\(|\bs{b}|=2\)，\(\bs{c}=\bs{a}+\bs{b}\)，且 \(\bs{c}\perp\bs{a}\)，则向量 \(\bs{a}\) 与 \(\bs{b}\) 的夹角为（\quad）

\choices
    {\(30^{\circ}\)}
    {\(60^{\circ}\)}
    {\(120^{\circ}\)}
    {\(150^{\circ}\)}
\begin{answer}
C
\end{answer}
\end{problem}

\begin{solution}
由 \(\bs{c}\perp\bs{a}\) 得 \((\bs{a}+\bs{b})\cdot\bs{a}=0\)，从而 \(\bs{a}\cdot\bs{b}=-|\bs{a}|^2=-1\). 若 \(\theta\) 是 \(\bs{a}\) 与 \(\bs{b}\) 的夹角，则 \(\cos\theta=\dfrac{\bs{a}\cdot\bs{b}}{|\bs{a}||\bs{b}|}=-\dfrac12\)，所以 \(\theta=120^{\circ}\)，故选 C.
\end{solution}

\begin{problem}
从原点向圆 \(x^2+y^2-12y+27=0\) 作两条切线，则这两条切线的夹角的大小为（\quad）

\choices
    {\(\frac{\pi}{6}\)}
    {\(\frac{\pi}{3}\)}
    {\(\frac{\pi}{2}\)}
    {\(\frac{2\pi}{3}\)}
\begin{answer}
B
\end{answer}
\end{problem}

\begin{solution}
将圆的方程化为 \(x^2+(y-6)^2=9\)，所以圆心为 \(C(0,6)\)，半径为 \(3\)，且 \(OC=6\). 设一条切线与圆相切于 \(T\)，则 \(CT\perp OT\). 在直角三角形 \(OCT\) 中，\(\sin\angle COT=\dfrac{CT}{OC}=\dfrac12\)，故 \(\angle COT=30^{\circ}\). 两条切线关于 \(OC\) 对称，夹角为 \(2\angle COT=60^{\circ}=\frac{\pi}{3}\)，故选 B.
\end{solution}

\begin{problem}
对任意的锐角 \(\alpha\)，\(\beta\)，下列不等关系中正确的是（\quad）

\choices
    {\(\sin(\alpha+\beta)>\sin\alpha+\sin\beta\)}
    {\(\sin(\alpha+\beta)>\cos\alpha+\cos\beta\)}
    {\(\cos(\alpha+\beta)<\sin\alpha+\sin\beta\)}
    {\(\cos(\alpha+\beta)<\cos\alpha+\cos\beta\)}
\begin{answer}
D
\end{answer}
\end{problem}

\begin{solution}
因为 \(\alpha\)，\(\beta\) 都是锐角，所以
\(\sin(\alpha+\beta)-\sin\alpha-\sin\beta
=\sin\alpha(\cos\beta-1)+\sin\beta(\cos\alpha-1)<0\)，故 A 不成立.

取 \(\alpha=\beta=30^{\circ}\)，则 \(\sin(\alpha+\beta)=\frac{\sqrt3}{2}<\sqrt3=\cos\alpha+\cos\beta\)，故 B 不成立. 取 \(\alpha=\beta=15^{\circ}\)，则 \(\cos(\alpha+\beta)=\frac{\sqrt3}{2}\)，而 \(\sin\alpha+\sin\beta=2\sin15^{\circ}=\frac{\sqrt6-\sqrt2}{2}<\frac{\sqrt3}{2}\)，故 C 也不成立.

另一方面，\(\cos\alpha>0\)，\(\cos\beta>0\)，且
\(\cos(\alpha+\beta)=\cos\alpha\cos\beta-\sin\alpha\sin\beta<\cos\alpha\cos\beta<\cos\alpha+\cos\beta\)，所以 D 正确.
\end{solution}

\begin{problem}
在正四面体 \(P-ABC\) 中，\(D\)，\(E\)，\(F\) 分别是 \(AB\)，\(BC\)，\(CA\) 的中点，下面四个结论中不成立的是（\quad）

\choices
    {\(BC\parallel\) 平面 \(PDF\)}
    {\(DF\perp\) 平面 \(PAE\)}
    {平面 \(PDF\)\(\perp\) 平面 \(ABC\)}
    {平面 \(PAE\)\(\perp\) 平面 \(ABC\)}
\begin{answer}
C
\end{answer}
\end{problem}

\begin{solution}
在正三角形 \(ABC\) 中，\(D\)，\(E\)，\(F\) 分别是 \(AB\)，\(BC\)，\(CA\) 的中点，所以 \(DF\parallel BC\). 又 \(P\notin\) 平面 \(ABC\)，而平面 \(PDF\) 与平面 \(ABC\) 的交线为 \(DF\)，故 \(C\notin\) 平面 \(PDF\). 因此 \(BC\parallel\) 平面 \(PDF\)，结论 A 成立.

设 \(O\) 是正三角形 \(ABC\) 的中心. 因为 \(AE\perp BC\)，且 \(DF\parallel BC\)，所以 \(DF\perp AE\). 又正四面体的高 \(PO\perp\) 平面 \(ABC\)，且 \(O\in AE\)，所以 \(PO\) 在平面 \(PAE\) 内. 令 \(Q=AE\cap DF\)，过 \(Q\) 作直线 \(q\parallel PO\)，则 \(q\subset\) 平面 \(PAE\)，并且 \(DF\perp q\). 直线 \(AE\) 与 \(q\) 在 \(Q\) 相交，故由直线与平面垂直的判定定理，\(DF\perp\) 平面 \(PAE\)，结论 B 成立. 又因为平面 \(PAE\) 内的直线 \(PO\perp\) 平面 \(ABC\)，所以平面 \(PAE\perp\) 平面 \(ABC\)，结论 D 成立.

验证结论 C 不成立. 取棱长为 \(2\)，建立直角坐标系，使 \(B=(-1,0,0)\)，\(C=(1,0,0)\)，\(A=(0,\sqrt3,0)\)，\(P=(0,\frac{\sqrt3}{3},\frac{2\sqrt6}{3})\). 则 \(D=(-\frac12,\frac{\sqrt3}{2},0)\)，\(F=(\frac12,\frac{\sqrt3}{2},0)\)，从而 \(\overrightarrow{DF}=(1,0,0)\)，\(\overrightarrow{DP}=(\frac12,-\frac{\sqrt3}{6},\frac{2\sqrt6}{3})\). 平面 \(PDF\) 的一个法向量为 \(\overrightarrow{DF}\times\overrightarrow{DP}=(0,-\frac{2\sqrt6}{3},-\frac{\sqrt3}{6})\)，它与平面 \(ABC\) 的法向量 \((0,0,1)\) 的内积不为 \(0\)，故两平面不垂直. 因此不成立的是 C.
\end{solution}

\begin{problem}
五个工程队承建某项工程的 \(5\) 个不同的子项目，每个工程队承建 \(1\) 项，其中甲工程队不能承建 \(1\) 号子项目，则不同的承建方案共有（\quad）

\choices
    {\(\mathrm C_4^1\mathrm C_4^4\)种}
    {\(\mathrm C_4^1\mathrm A_4^4\)种}
    {\(\mathrm C_4^4\)种}
    {\(\mathrm A_4^4\)种}
\begin{answer}
B
\end{answer}
\end{problem}

\begin{solution}
若不考虑限制条件，把 \(5\) 个工程队分配给 \(5\) 个不同子项目，共有 \(5!=120\) 种方案. 甲工程队承建 \(1\) 号子项目时，其他 \(4\) 个工程队分配给其余 \(4\) 个子项目，有 \(4!=24\) 种方案. 因此符合条件的方案数为 \(5!-4!=120-24=96\). 又 \(\mathrm C_4^1\mathrm A_4^4=4\times24=96\)，故选 B.
\end{solution}

\section{填空题}
\begin{problem}
抛物线 \(y^2=4x\) 的准线方程是\fillinblank{}，焦点坐标是\fillinblank{}.
\begin{answer}
\(x=-1\)，\((1,0)\)
\end{answer}
\end{problem}

\begin{solution}
抛物线的标准方程为 \(y^2=4px\). 与 \(y^2=4x\) 比较得 \(p=1\)，所以准线方程为 \(x=-p=-1\)，焦点为 \((p,0)=(1,0)\).
\end{solution}

\begin{problem}
\(\left(x-\frac{1}{x}\right)^6\) 的展开式中的常数项是\fillinblank{}.（用数字作答）
\begin{answer}
\(-20\)
\end{answer}
\end{problem}

\begin{solution}
展开式的通项为 \(\mathrm C_6^k x^{6-k}\left(-x^{-1}\right)^k=\mathrm C_6^k(-1)^k x^{6-2k}\). 要得到常数项，必须有 \(6-2k=0\)，即 \(k=3\)，所以常数项为 \(\mathrm C_6^3(-1)^3=-20\).
\end{solution}

\begin{problem}
函数 \(f(x)=\sqrt{x+1}+\frac{1}{2-x}\) 的定义域为\fillinblank{}.
\begin{answer}
\([-1,2)\cup(2,+\infty)\)
\end{answer}
\end{problem}

\begin{solution}
根式有意义要求 \(x+1\geqslant0\)，即 \(x\geqslant-1\)；分式有意义要求 \(2-x\neq0\)，即 \(x\neq2\). 合并得定义域为 \([-1,2)\cup(2,+\infty)\).
\end{solution}

\begin{problem}
在 \(\triangle ABC\) 中，\(AC=\sqrt{3}\)，\(\angle A=45^{\circ}\)，\(\angle C=75^{\circ}\)，则 \(BC\) 的长为\fillinblank{}.
\begin{answer}
\(\sqrt{2}\)
\end{answer}
\end{problem}

\begin{solution}
由三角形内角和得 \(\angle B=180^{\circ}-45^{\circ}-75^{\circ}=60^{\circ}\). 由正弦定理，\(\dfrac{BC}{\sin45^{\circ}}=\dfrac{AC}{\sin60^{\circ}}\)，所以 \(BC=\sqrt3\cdot\dfrac{\sqrt2}{2}\div\dfrac{\sqrt3}{2}=\sqrt2\).
\end{solution}

\begin{problem}
对于函数 \(f(x)\) 定义域中任意的 \(x_1\)，\(x_2\)（\(x_1\neq x_2\)），有如下结论：

\begin{circlelist}
\item \(f(x_1+x_2)=f(x_1)\cdot f(x_2)\)；
\item \(f(x_1\cdot x_2)=f(x_1)+f(x_2)\)；
\item \(\dfrac{f(x_1)-f(x_2)}{x_1-x_2}>0\)；
\item \(f\left(\dfrac{x_1+x_2}{2}\right)<\dfrac{f(x_1)+f(x_2)}{2}\).
\end{circlelist}

当 \(f(x)=\lg x\) 时，上述结论中正确结论的序号是\fillinblank{}.
\begin{answer}
\(\circled{2}\circled{3}\)
\end{answer}
\end{problem}

\begin{solution}
因为 \(f(x)=\lg x\) 的定义域为 \((0,+\infty)\). 结论 \(\circled{1}\) 一般不成立，例如取 \(x_1=1\)，\(x_2=2\) 时，左边为 \(\lg3>0\)，右边为 \(\lg1\cdot\lg2=0\). 结论 \(\circled{2}\) 成立，因为 \(\lg(x_1x_2)=\lg x_1+\lg x_2\).

函数 \(f(x)=\lg x\) 在 \((0,+\infty)\) 上单调递增，所以当 \(x_1\neq x_2\) 时，\(\dfrac{\lg x_1-\lg x_2}{x_1-x_2}>0\)，结论 \(\circled{3}\) 成立. 又因 \(x_1\)，\(x_2\) 为正数且 \(x_1\neq x_2\)，由算术—几何平均不等式得 \(\dfrac{x_1+x_2}{2}>\sqrt{x_1x_2}\)，从而 \(\lg\left(\dfrac{x_1+x_2}{2}\right)>\lg\sqrt{x_1x_2}=\dfrac{1}{2}\lg(x_1x_2)=\dfrac{\lg x_1+\lg x_2}{2}\)，结论 \(\circled{4}\) 的不等号方向相反. 因此正确结论的序号是 \(\circled{2}\)、\(\circled{3}\).
\end{solution}

\begin{problem}
已知 \(n\) 次多项式 \(P_n(x)=a_0x^n+a_1x^{n-1}+\cdots+a_{n-1}x+a_n\).

如果在一种算法中，计算 \(x_0^k\)（\(k=2\)，\(3\)，\(4\)，\(\cdots\)，\(n\)）的值需要 \(k-1\) 次乘法，计算 \(P_3(x_0)\) 的值共需要 \(9\) 次运算（\(6\) 次乘法，\(3\) 次加法），那么计算 \(P_{10}(x_0)\) 的值共需要\fillinblank{} 次运算.

下面给出一种减少运算次数的算法：\(P_0(x)=a_0\)，\(P_{k+1}(x)=xP_k(x)+a_{k+1}\)（\(k=0\)，\(1\)，\(2\)，\(\cdots\)，\(n-1\)）. 利用该算法，计算 \(P_3(x_0)\) 的值共需要 \(6\) 次运算，计算 \(P_{10}(x_0)\) 的值共需要\fillinblank{} 次运算.
\begin{answer}
\(65\)，\(20\)
\end{answer}
\end{problem}

\begin{solution}
按第一种算法，先计算 \(x_0^2\)，\(\ldots\)，\(x_0^{10}\)，所需乘法次数为 \(1+2+\cdots+9=45\). 随后计算 \(a_0x_0^{10}\)，\(a_1x_0^9\)，\(\ldots\)，\(a_9x_0\)，还需 \(10\) 次乘法；把这 \(10\) 项与 \(a_{10}\) 相加，需 \(10\) 次加法. 因此总运算次数为 \(45+10+10=65\).

按递推算法，从 \(P_0(x_0)\) 逐步得到 \(P_1(x_0)\)，\(\ldots\)，\(P_{10}(x_0)\)，每一步 \(P_{k+1}(x_0)=x_0P_k(x_0)+a_{k+1}\) 都需 \(1\) 次乘法和 \(1\) 次加法，共有 \(10\) 步，所以总运算次数为 \(10+10=20\).
\end{solution}

\section{解答题}
\begin{problem}
已知 \(\tan\frac{\alpha}{2}=2\)，求：

\begin{enumerate}
\item \(\tan\left(\alpha+\frac{\pi}{4}\right)\) 的值；
\item \(\dfrac{6\sin\alpha+\cos\alpha}{3\sin\alpha-2\cos\alpha}\) 的值.
\end{enumerate}
\end{problem}

\begin{solution}
由半角正切公式，\(\sin\alpha=\dfrac{2\tan\frac{\alpha}{2}}{1+\tan^2\frac{\alpha}{2}}=\dfrac45\)，\(\cos\alpha=\dfrac{1-\tan^2\frac{\alpha}{2}}{1+\tan^2\frac{\alpha}{2}}=-\dfrac35\)，所以 \(\tan\alpha=\dfrac{\sin\alpha}{\cos\alpha}=-\dfrac43\).

\begin{enumerate}
\item \(\tan\left(\alpha+\frac{\pi}{4}\right)=\dfrac{\tan\alpha+1}{1-\tan\alpha}=\dfrac{-\frac43+1}{1+\frac43}=-\dfrac17\).
\item \(\dfrac{6\sin\alpha+\cos\alpha}{3\sin\alpha-2\cos\alpha}=\dfrac{6\cdot\frac45-\frac35}{3\cdot\frac45-2\left(-\frac35\right)}=\dfrac{21}{18}=\dfrac76\).
\end{enumerate}
\end{solution}

\begin{problem}
如图，在直三棱柱 \(ABC-A_1B_1C_1\) 中，\(AC=3\)，\(BC=4\)，\(AB=5\)，\(AA_1=4\)，点 \(D\) 是 \(AB\) 的中点.

\begin{enumerate}
\item 求证：\(AC\perp BC_1\)；
\item 求证：\(AC_1\parallel\) 平面 \(CDB_1\)；
\item 求异面直线 \(AC_1\) 与 \(B_1C\) 所成角的余弦值.
\end{enumerate}

\bitmapfigure[max width=0.45\linewidth]{img/2005/beijing_liberal_autumn/q16_fig1.png}
\end{problem}

\begin{solution}
建立空间直角坐标系，取 \(C\) 为原点，令 \(A=(3,0,0)\)，\(B=(0,4,0)\)，并令 \(z\) 轴垂直于底面 \(ABC\). 则 \(C_1=(0,0,4)\)，\(B_1=(0,4,4)\)，\(D=(\frac32,2,0)\).

\begin{enumerate}
\item \(\overrightarrow{AC}=(-3,0,0)\)，\(\overrightarrow{BC_1}=(0,-4,4)\)，所以 \(\overrightarrow{AC}\cdot\overrightarrow{BC_1}=0\)，故 \(AC\perp BC_1\).
\item \(\overrightarrow{AC_1}=(-3,0,4)\)，\(\overrightarrow{CD}=(\frac32,2,0)\)，\(\overrightarrow{CB_1}=(0,4,4)\)，且 \(\overrightarrow{AC_1}=-2\overrightarrow{CD}+\overrightarrow{CB_1}\). 因此直线 \(AC_1\) 的方向向量是平面 \(CDB_1\) 内两个相交方向向量的线性组合. 另外，若点 \(A\) 在平面 \(CDB_1\) 内，则由 \(\overrightarrow{CA}=(3,0,0)=u\overrightarrow{CD}+v\overrightarrow{CB_1}\) 的竖直坐标得 \(v=0\)，进而 \(\overrightarrow{CA}=u\overrightarrow{CD}\)，但 \((3,0,0)\) 不是 \((\frac32,2,0)\) 的倍数，矛盾. 所以 \(A\notin\) 平面 \(CDB_1\)，故 \(AC_1\parallel\) 平面 \(CDB_1\).
\item 取 \(\overrightarrow{B_1C}=(0,-4,-4)\)，则 \(|\overrightarrow{AC_1}|=5\)，\(|\overrightarrow{B_1C}|=4\sqrt2\)，且 \(\left|\overrightarrow{AC_1}\cdot\overrightarrow{B_1C}\right|=16\). 所以异面直线 \(AC_1\) 与 \(B_1C\) 所成角的余弦值为 \(\dfrac{16}{5\cdot4\sqrt2}=\dfrac{2\sqrt2}{5}\).
\end{enumerate}
\end{solution}

\begin{problem}
数列 \(\{a_n\}\) 的前 \(n\) 项和为 \(S_n\)，且 \(a_1=1\)，\(a_{n+1}=\frac{1}{3}S_n\)，\(n=1\)，\(2\)，\(3\)，\(\cdots\)，求：

\begin{enumerate}
\item \(a_2\)，\(a_3\)，\(a_4\) 的值及数列 \(\{a_n\}\) 的通项公式；
\item \(a_2+a_4+a_6+\cdots+a_{2n}\) 的值.
\end{enumerate}
\end{problem}

\begin{solution}
\begin{enumerate}
\item \(a_2=\frac13S_1=\frac13\). 又 \(S_n=S_{n-1}+a_n=S_{n-1}+\frac13S_{n-1}=\frac43S_{n-1}\)（\(n\geqslant2\)），所以 \(a_{n+1}=\frac43a_n\)（\(n\geqslant2\)）. 因此 \(a_3=\frac49\)，\(a_4=\frac{16}{27}\)，且当 \(n\geqslant2\) 时，\(a_n=\frac13\left(\frac43\right)^{n-2}\).
\item \(\displaystyle a_2+a_4+\cdots+a_{2n}=\frac13\left[1+\left(\frac43\right)^2+\cdots+\left(\frac43\right)^{2n-2}\right]=\frac13\cdot\frac{\left(\frac43\right)^{2n}-1}{\left(\frac43\right)^2-1}=\frac37\left[\left(\frac43\right)^{2n}-1\right]\).
\end{enumerate}
\end{solution}

\begin{problem}
甲，乙两人各进行 \(3\) 次射击，甲每次击中目标的概率为 \(\frac{1}{2}\)，乙每次击中目标的概率为 \(\frac{2}{3}\).

\begin{enumerate}
\item 甲恰好击中目标 \(2\) 次的概率；
\item 乙至少击中目标 \(2\) 次的概率；
\item 乙恰好比甲多击中目标 \(2\) 次的概率.
\end{enumerate}
\end{problem}

\begin{solution}
\begin{enumerate}
\item 甲恰好击中 \(2\) 次的概率为 \(\mathrm C_3^2\left(\frac12\right)^2\left(\frac12\right)=\frac38\).
\item 乙恰好击中 \(2\) 次的概率为 \(\mathrm C_3^2\left(\frac23\right)^2\frac13=\frac49\)，乙击中 \(3\) 次的概率为 \(\left(\frac23\right)^3=\frac{8}{27}\)，所以乙至少击中 \(2\) 次的概率为 \(\frac49+\frac{8}{27}=\frac{20}{27}\).
\item 设甲、乙击中目标的次数分别为 \(X\)，\(Y\). 条件 \(Y-X=2\) 只有两种情况：\(X=0\)，\(Y=2\) 或 \(X=1\)，\(Y=3\). 因此所求概率为 \(\left(\frac12\right)^3\mathrm C_3^2\left(\frac23\right)^2\frac13+\mathrm C_3^1\frac12\left(\frac12\right)^2\left(\frac23\right)^3=\frac1{18}+\frac19=\frac16\).
\end{enumerate}
\end{solution}

\begin{problem}
已知函数 \(f(x)=-x^3+3x^2+9x+a\).

\begin{enumerate}
\item 求 \(f(x)\) 的单调减区间；
\item 若 \(f(x)\) 在区间 \([-2,2]\) 上的最大值为 \(20\)，求它在该区间上的最小值.
\end{enumerate}
\end{problem}

\begin{solution}
\begin{enumerate}
\item \(f'(x)=-3x^2+6x+9=-3(x-3)(x+1)\). 当 \(x<-1\) 或 \(x>3\) 时，\(f'(x)<0\)，所以 \(f(x)\) 的单调减区间为 \((-\infty,-1)\) 和 \((3,+\infty)\).
\item 在 \([-2,2]\) 上，\(f(x)\) 在 \([-2,-1]\) 上递减，在 \([-1,2]\) 上递增，因此最大值取在端点. 又 \(f(-2)=a+2\)，\(f(2)=a+22\)，所以最大值为 \(a+22=20\)，得 \(a=-2\). 最小值取在 \(x=-1\)，为 \(f(-1)=1+3-9-2=-7\).
\end{enumerate}
\end{solution}

\begin{problem}
如图，直线 \(l_1:y=kx\)（\(k>0\)）与直线 \(l_2:y=-kx\) 之间的阴影区域（不含边界）记为 \(W\)，其左半部分记为 \(W_1\)，右半部分记为 \(W_2\).

\begin{enumerate}
\item 分别用不等式组表示 \(W_1\) 和 \(W_2\)；
\item 若区域 \(W\) 中的动点 \(P(x,y)\) 到 \(l_1\)，\(l_2\) 的距离之积等于 \(d^2\)，求点 \(P\) 的轨迹 \(C\) 的方程；
\item 设不过原点 \(O\) 的直线 \(l\) 与（2）中的曲线 \(C\) 相交于 \(M_1\)，\(M_2\) 两点，且与 \(l_1\)，\(l_2\) 分别交于 \(M_3\)，\(M_4\) 两点. 求证：\(\triangle OM_1M_2\) 的重心与 \(\triangle OM_3M_4\) 的重心重合.
\end{enumerate}

\bitmapfigure[max width=0.42\linewidth]{img/2005/beijing_liberal_autumn/q20_fig1.png}
\end{problem}

\begin{solution}
\begin{enumerate}
\item 当 \(x<0\) 时，\(kx<y<-kx\)，所以 \(W_1=\{(x,y)\mid kx<y<-kx,\ x<0\}\)；当 \(x>0\) 时，\(-kx<y<kx\)，所以 \(W_2=\{(x,y)\mid-kx<y<kx,\ x>0\}\).
\item 两直线的方程分别为 \(kx-y=0\) 和 \(kx+y=0\)，所以点 \(P(x,y)\) 到它们的距离之积为 \(\dfrac{|kx-y||kx+y|}{k^2+1}\). 在区域 \(W\) 内有 \(y^2<k^2x^2\)，从而 \(k^2x^2-y^2=(k^2+1)d^2\)，即轨迹 \(C\) 的方程为 \(k^2x^2-y^2-(k^2+1)d^2=0\).
\item 先设 \(l\) 不是竖直直线，写成 \(y=mx+b\). 因 \(l\) 不过原点，\(b\neq0\)；又因 \(l\) 与双曲线 \(C\) 有两个交点，\(m\neq\pm k\). 把 \(y=mx+b\) 代入 \(C\)，交点横坐标 \(x_1\)，\(x_2\) 是方程
\( (k^2-m^2)x^2-2mbx-\left[b^2+(k^2+1)d^2\right]=0 \)
的两个根，所以 \(x_1+x_2=\dfrac{2mb}{k^2-m^2}\).

直线 \(l\) 与 \(l_1\)，\(l_2\) 的交点横坐标分别为 \(x_3=\dfrac{b}{k-m}\)，\(x_4=-\dfrac{b}{m+k}\)，从而 \(x_3+x_4=\dfrac{2mb}{k^2-m^2}=x_1+x_2\). 由于四点都在直线 \(l\) 上，\(y_i=mx_i+b\)，故 \(y_1+y_2=m(x_1+x_2)+2b=m(x_3+x_4)+2b=y_3+y_4\). 于是两个三角形的三个顶点坐标和相同，且都含有顶点 \(O\)，所以它们的重心坐标分别为 \(\left(\dfrac{x_1+x_2}{3},\dfrac{y_1+y_2}{3}\right)\) 和 \(\left(\dfrac{x_3+x_4}{3},\dfrac{y_3+y_4}{3}\right)\)，二者相同.

若 \(l\) 是竖直直线 \(x=c\)，因 \(l\) 不过原点有 \(c\neq0\). 曲线 \(C\) 与它的两个交点的纵坐标互为相反数，故 \(y_1+y_2=0\)；而 \(M_3=(c,kc)\)，\(M_4=(c,-kc)\)，也有 \(y_3+y_4=0\)，并且 \(x_1+x_2=x_3+x_4=2c\). 因此这时两个三角形的重心也重合.
\end{enumerate}
\end{solution}
